% FILE: manuscripts/2026-thesis/project-description/figures-tex/figure-style.tex
% ==================================================================================================
% Shared layer of the figures of the thesis proposal: canvas, colours, type scale, TikZ styles and
% drawing macros. A figure defines \RootPath and \CanvasHeight, then inputs this file.
%
% The visual grammar follows the PINTS 2026 graphical abstract (manuscripts/2026-11_pints/figure.tex):
% the same poster colour layers, the same camera and the same vector styles, so that a reader of
% both artefacts meets one code. The PINTS figure declares its styles inside its own document, so
% the styles below re-declare them at the scale of an A4 page; a style file shared by both
% artefacts remains to be extracted.

\usepackage{geometry}
\input{\RootPath/common/preambles/user}
\usepackage{tikz-3dplot}
\usetikzlibrary{calc,arrows.meta,shadows.blur}
\usepackage[nottdefault]{sourcecodepro}
% Every label is set through \sffamily, which Source Sans Pro now serves, so that the lettering of the
% figures matches the text of the proposal (full-text-tex/proposal-style.tex).
\usepackage[semibold]{sourcesanspro}
\usepackage{fontawesome5}
\newcommand{\KickerFont}{\fontfamily{SourceCodePro-TLF}\selectfont}

% ---- Canvas --------------------------------------------------------------------------------------
% The width is the text width of the A4 proposal, so that each figure prints at true scale. Scholia
% configures an A4 page when the common preamble loads; the canvas is applied after that page.
\newcommand{\CanvasWidth}{16}
\geometry{paperwidth=\CanvasWidth cm,paperheight=\CanvasHeight cm,margin=0pt}
% Scholia stretches every leading by 1.2 for running prose; the type scale below sets each leading by
% role, so the stretch is reset here.
\setstretch{1}
\pagestyle{empty}
\setlength{\parindent}{0pt}

% ---- Colours -------------------------------------------------------------------------------------
% The poster's value, scheme and role layers are the single source of every colour token. The order
% follows the poster's parts/colors.tex. That file cannot be loaded from here, because its inputs
% are relative to the poster folder. The scheme is the petrol-guide variant used by the PINTS figure.
\newcommand{\PosterParts}{\RootPath/manuscripts/2026-08-ccn-extended-abstract/3-poster/parts}
\input{\PosterParts/colors-values}
\newif\ifquestioninverted\questioninvertedtrue
\input{\PosterParts/../schemes/petrol-guide}
\input{\PosterParts/colors-roles}
% The poster's macro layer supplies \taperarrow, the geometry of the connectors.
\input{\PosterParts/macros}

% Roles of the proposal figures, bound to existing semantic tokens. One concept, one colour, across
% the three figures:
%   - the two factors of a task and the two components of a mixture share the component colours;
%   - the interaction component and the emergent direction share one colour, since each is the part
%     of a representation that neither factor or component produces alone;
%   - the chrome marks the elements of an explanation (boxes, criteria), and never a concept.
\colorlet{factorone}{comp1}
\colorlet{factortwo}{comp2}
\colorlet{conjunction}{outplane}
\colorlet{meanpart}{borderfunc}
\colorlet{mixtext}{mixresponse!80!black}
\colorlet{figchrome}{chromecore}
\colorlet{pageground}{surface}

% ---- Type scale ----------------------------------------------------------------------------------
% THE ONLY PLACE A POINT SIZE IS CHOSEN in these figures. Each figure prints at true scale, so a
% point size is a physical size. Two floors:
%   WORDS (headings, notes, legends, tags): 7 pt;   MARKS inside a diagram (ticks, axis names): 6 pt.
\newcommand{\HeadSize}{\fontsize{8.5}{9.8}\selectfont}    % card headings, panel letters
\newcommand{\BodySize}{\fontsize{7.5}{8.6}\selectfont}    % words about a diagram (compact leading)
\newcommand{\ObjectSize}{\fontsize{8}{8}\selectfont}      % labels of the drawn objects (flat)
\newcommand{\TickSize}{\fontsize{6.5}{7.5}\selectfont}    % subordinate marks -- above the 6 pt floor
\newcommand{\KickerSize}{\fontsize{8}{9.2}\selectfont}    % operation strips, in the kicker face
\newcommand{\IconSize}{\fontsize{7}{7}\selectfont}        % icon glyphs in a disc: exempt from floors
\DeclareMathSizes{8.5}{9.5}{7}{5}
\DeclareMathSizes{8}{9}{6.5}{5}
\DeclareMathSizes{7.5}{8.4}{6}{5}
\DeclareMathSizes{6.5}{7.3}{5.5}{5}

% ---- Styles --------------------------------------------------------------------------------------
\tikzset{
  % Cards: one panel per function. A dashed rim marks a claim labelled in the proposal as a
  % hypothesis or as prospective; a solid rim marks an established claim.
  card/.style={fill=surface,draw=bordersubtle,line width=0.8pt,rounded corners=6pt},
  card open/.style={fill=surface,draw=borderfunc,line width=0.8pt,rounded corners=6pt,
    dash pattern=on 3pt off 2pt},
  % A legend that applies to every panel: a ground of its own, a functional rim and no shadow, so
  % that it reads apart from the lettered panels.
  legend card/.style={fill=posterbg,draw=borderfunc,line width=0.6pt,rounded corners=6pt},
  card letter/.style={circle,fill=figchrome,text=white,inner sep=0pt,minimum size=0.40cm,
    font=\sffamily\bfseries\HeadSize},
  heading/.style={anchor=base west,inner sep=0pt,font=\sffamily\bfseries\HeadSize,
    text=textprimary},
  subheading/.style={anchor=base west,inner sep=0pt,font=\sffamily\BodySize,text=textsecondary},
  status tag/.style={anchor=north east,font=\sffamily\BodySize,text=texttertiary,
    draw=borderfunc,rounded corners=2pt,inner xsep=2.5pt,inner ysep=1.5pt,
    dash pattern=on 1.5pt off 1pt},
  body/.style={font=\sffamily\BodySize,text=textsecondary,align=center},
  keyword/.style={font=\sffamily\bfseries\BodySize,align=center},
  object/.style={font=\ObjectSize},
  tick/.style={font=\sffamily\TickSize,text=texttertiary},
  % Vectors and planes of activity space (after the PINTS figure, at 0.8 of its widths).
  axis3d/.style={-{Stealth[length=1.6mm,width=1.2mm]},draw=borderfunc,line width=0.6pt},
  plane/.style={fill=compplane,fill opacity=0.35,draw=compplane!80!black,draw opacity=0.5,
    line width=0.5pt},
  guide/.style={draw=texttertiary!80,line width=0.55pt,dash pattern=on 1.6pt off 1.2pt},
  edge/.style={draw=bordersubtle!70!borderfunc,line width=0.6pt},
  comp one/.style={-{Stealth[length=2.0mm,width=1.5mm]},draw=factorone,line width=1.3pt},
  comp two/.style={-{Stealth[length=2.0mm,width=1.5mm]},draw=factortwo,line width=1.3pt},
  mixture/.style={-{Stealth[length=2.3mm,width=1.7mm]},draw=mixresponse,line width=1.7pt},
  reference/.style={-{Stealth[length=2.0mm,width=1.5mm]},draw=textsecondary,line width=1.1pt,
    dash pattern=on 2.6pt off 1.5pt},
  orthogonal/.style={-{Stealth[length=1.9mm,width=1.4mm]},draw=conjunction,line width=1.3pt},
  flow/.style={-{Stealth[length=2.0mm,width=1.6mm]},draw=textsecondary,line width=0.9pt},
  leader/.style={line width=0.5pt,shorten >=1.5pt,shorten <=1.5pt},
  coordinate dot/.style={circle,fill=textsecondary,inner sep=0.9pt},
  % Conditions of a task, and their target outputs: a filled disc for one output value, an open
  % disc for the other. The colour names the mapping whose output the disc shows.
  cond on/.style={circle,fill=#1,draw=#1,line width=0.8pt,inner sep=0pt,minimum size=2.1mm},
  cond on/.default=textprimary,
  cond off/.style={circle,fill=white,draw=#1,line width=0.8pt,inner sep=0pt,minimum size=2.1mm},
  cond off/.default=textprimary,
  readout plane/.style={fill=#1,fill opacity=0.20,draw=#1,draw opacity=0.75,line width=0.6pt},
  % Elements of an explanation.
  element/.style={rounded corners=4pt,draw=figchrome,line width=0.8pt,fill=figchrome!7,
    align=center,font=\sffamily\BodySize,text=textprimary,inner sep=2.5pt},
  determinant/.style={rounded corners=4pt,draw=borderfunc,line width=0.6pt,fill=surface,
    align=center,font=\sffamily\BodySize,text=textprimary,inner xsep=4pt,inner ysep=2pt},
  criterion/.style={font=\sffamily\slshape\BodySize,text=figchrome,align=center},
  relation/.style={font=\sffamily\bfseries\BodySize,text=figchrome},
  pairing arrow/.style={-{Stealth[length=2.4mm,width=2.0mm]},draw=figchrome,line width=1.4pt},
  % Operation strips of the recurrent network (the poster's stage strips).
  operation/.style={rounded corners=5pt,minimum height=0.56cm,font=\KickerFont\KickerSize},
  neuron/.style={circle,draw=textsecondary,fill=surface,line width=0.5pt,minimum size=0.12cm,
    inner sep=0pt},
  verdict yes/.style={circle,fill=conjunction,text=white,inner sep=0pt,minimum size=0.36cm,
    font=\IconSize},
  verdict no/.style={circle,fill=borderfunc,text=white,inner sep=0pt,minimum size=0.36cm,
    font=\IconSize},
}

% One orthographic camera for every 3-D plot, as in the PINTS figure: the third axis points up, the
% first towards the lower left and the second towards the lower right.
\tdplotsetmaincoords{60}{122}

% ---- Canvas environment --------------------------------------------------------------------------
% Coordinates in cm from the page's south-west corner.
\newenvironment{FigureCanvas}{%
  \thispagestyle{empty}%
  \begin{tikzpicture}[remember picture,overlay,line cap=round,line join=round]
  \begin{scope}[shift={(current page.south west)},x=1cm,y=1cm]
  \fill[pageground] (0,0) rectangle (\CanvasWidth,\CanvasHeight);
}{%
  \end{scope}
  \end{tikzpicture}%
  \null}

% ---- Cards ---------------------------------------------------------------------------------------
% \SoftShadow{<x1>}{<y1>}{<x2>}{<y2>} : a soft drop under the rounded rectangle (x1,y1)-(x2,y2), drawn
% before it, as stacked low-opacity layers (the PINTS figure's shadow, at a fixed depth).
\newcommand{\SoftShadow}[4]{%
  \foreach \g in {1,...,5}{%
    \pgfmathsetlengthmacro\ShadowCorner{6pt+0.22mm*\g}%
    \fill[black,opacity=0.018,rounded corners=\ShadowCorner]
      ({#1-0.02*\g},{#2-0.07-0.02*\g}) rectangle ({#3+0.02*\g},{#4-0.07+0.02*\g});
  }%
}
% Offsets of the card header from the card's top-left corner (cm): the letter disc, then the
% heading and subheading baselines. Every card reads them, so all headers align.
\newcommand{\CardLetterX}{0.34}
\newcommand{\CardLetterY}{0.34}
\newcommand{\CardTextX}{0.66}
\newcommand{\CardHeadY}{0.45}
\newcommand{\CardSubY}{0.80}
% \CardHeader{<x1>}{<y2>}{<letter>}{<heading>}{<subheading>}
\newcommand{\CardHeader}[5]{%
  \node[card letter] at ({#1+\CardLetterX},{#2-\CardLetterY}) {#3};
  \node[heading] at ({#1+\CardTextX},{#2-\CardHeadY}) {#4};
  \node[subheading] at ({#1+\CardTextX},{#2-\CardSubY}) {#5};
}
% \Card{<x1>}{<y1>}{<x2>}{<y2>}{<letter>}{<heading>}{<subheading>} : an established panel.
\newcommand{\Card}[7]{%
  \SoftShadow{#1}{#2}{#3}{#4}
  \fill[card] (#1,#2) rectangle (#3,#4);
  \draw[card] (#1,#2) rectangle (#3,#4);
  \CardHeader{#1}{#4}{#5}{#6}{#7}
}
% \OpenCard{<x1>}{<y1>}{<x2>}{<y2>}{<letter>}{<heading>}{<subheading>}{<status>} : a panel whose
% claim is a hypothesis or prospective, with its status named in the top-right corner.
\newcommand{\OpenCard}[8]{%
  \fill[card] (#1,#2) rectangle (#3,#4);
  \draw[card open] (#1,#2) rectangle (#3,#4);
  \CardHeader{#1}{#4}{#5}{#6}{#7}
  \node[status tag] at ({#3-0.15},{#4-0.15}) {#8};
}

% ---- Activity space ------------------------------------------------------------------------------
% \NeuronAxes{<x>}{<y>}{<z>} : three axes of activity space from the origin, each labelled beyond its
% tip, where no point of a representation lies. Called inside a tdplot_main_coords scope.
\newcommand{\NeuronAxes}[3]{%
  \draw[axis3d] (0,0,0) -- (#1,0,0) node[tick,anchor=north east,inner sep=0.5pt] {neuron 1};
  \draw[axis3d] (0,0,0) -- (0,#2,0) node[tick,anchor=north west,inner sep=0.5pt] {neuron 2};
  \draw[axis3d] (0,0,0) -- (0,0,#3) node[tick,anchor=south,inner sep=1pt] {neuron 3};
}
% \BareAxes{<x>}{<y>}{<z>} : the same axes, unlabelled, for a small repeated plot.
\newcommand{\BareAxes}[3]{%
  \draw[axis3d] (0,0,0) -- (#1,0,0);
  \draw[axis3d] (0,0,0) -- (0,#2,0);
  \draw[axis3d] (0,0,0) -- (0,0,#3);
}

% ---- Conditions of a task with two binary factors ------------------------------------------------
% The decision value d runs along x and the context value c along y, both in {-1, +1}. Every drawing
% of the conditions (the task square, the sign tiles) uses this placement, so a condition keeps one
% position across figures.
%
% \TaskGlyph{<cx>}{<cy>}{<half side>} : the four conditions, each a disc showing its target output
% under the exclusive-or (filled when d and c differ).
\newcommand{\TaskGlyph}[3]{%
  \draw[guide] ({#1-#3},{#2-#3}) rectangle ({#1+#3},{#2+#3});
  \foreach \d in {-1,1}{%
    \foreach \c in {-1,1}{%
      \pgfmathtruncatemacro\Differ{\d*\c<0}%
      \ifnum\Differ=1
        \node[cond on] at ({#1+\d*#3},{#2+\c*#3}) {};
      \else
        \node[cond off] at ({#1+\d*#3},{#2+\c*#3}) {};
      \fi
    }%
  }%
}
% \ConditionSquare{<cx>}{<cy>}{<half side>} : the task glyph with its two factor axes.
\newcommand{\ConditionSquare}[3]{%
  \TaskGlyph{#1}{#2}{#3}
  \draw[-{Stealth[length=1.6mm,width=1.3mm]},draw=factorone,line width=0.8pt]
    ({#1-#3-0.15},{#2-#3-0.38}) -- ({#1+#3+0.15},{#2-#3-0.38})
    node[midway,below,inner sep=1.5pt,font=\sffamily\BodySize,text=factorone] {decision};
  \draw[-{Stealth[length=1.6mm,width=1.3mm]},draw=factortwo,line width=0.8pt]
    ({#1-#3-0.38},{#2-#3-0.15}) -- ({#1-#3-0.38},{#2+#3+0.15})
    node[midway,above,sloped,inner sep=1.5pt,font=\sffamily\BodySize,text=factortwo] {context};
}
% ---- Connector -----------------------------------------------------------------------------------
% \Connector{<tail>}{<head>} : the poster's connector (\connarrow), at the widths of the PINTS figure
% scaled for runs of about 1 cm. The options are passed literally, so that pgfkeys splits the list.
\newcommand{\Connector}[2]{%
  \taperarrow[spinefill!45][draw=spineedge,line width=0.7pt]%
    {#1}{#2}{0pt}{0.6pt}{3.0pt}{5.5pt}{7.2pt}}
